\begin{lemma}[Defective-chain structure, part (c)]
Under the hypotheses of parts (a)--(b), direct the full path $P$ from
$u\in V(T)$ toward its outside end.  In this path order there is a first
vertex $v_1\notin V(T)$ and, strictly after that first departure, a last
vertex $v_2\in V(T)$.
\end{lemma}
\begin{proof}
The path starts at the tree vertex $u$ and ends at the outside vertex $z$
by \noderef{ScheideDefectiveChainPathExists}, so the finite ordered list of
vertex-nodes on the incidence path has a first member outside $V(T)$; call
it $v_1$.  The corresponding incidence index is the first-departure index.

By \noderef{ScheideGlobalMaximumBoundaryFacts}, the defective color
$\alpha$ has at least three distinct boundary edge identities, whereas no
$\gamma$-edge leaves $V(T)$.  By
\noderef{ScheideDefectiveChainBoundaryEdges}, all of those
$\alpha$-boundary edge-nodes occur on $P$.  Starting inside and crossing
the boundary at least three times before ending outside forces a return to
$V(T)$ after the first departure.  Hence there is a tree vertex after
$v_1$.  Among the finitely many such later vertex-nodes choose the last
one and call it $v_2$.  The first and last choices give exactly the index
conditions in \noderef{ScheideDepartureReturn}.  Notice that the indices
are incidence-walk indices: boundary edge-nodes occur between consecutive
multigraph vertex-nodes, so no crossing is lost by parallelism.
\end{proof}
